%% -------------------------------------------------------------------
%% Hierarchical Transformer Architectures for ICAO ADREP Aviation
%% Incident Classification
%% Team 30, 11-785 (Spring 2026)
%% -------------------------------------------------------------------
%% v2 draft — addresses TA feedback from the midterm:
%%   (TA1) author emails, Aviation-BERT integration, focal-loss gamma
%%         sweep, learnable gating for compliance/accuracy trade-off.
%%   (TA2) fix broken `jing2023bert` citation, expand Related Work with
%%         proper citations, reconcile Figure 1 with long-tail claim.
%% -------------------------------------------------------------------

\documentclass{article}

% Use [preprint] to suppress the "33rd Conference on NeurIPS 2019" footer;
% this is a course paper, not a NeurIPS submission.
\usepackage[preprint]{neurips_2019}

\usepackage[utf8]{inputenc}
\usepackage[T1]{fontenc}
\usepackage{hyperref}
\usepackage{url}
\usepackage{booktabs}
\usepackage{amsmath}
\usepackage{amsfonts}
\usepackage{nicefrac}
\usepackage{microtype}
\usepackage{graphicx}
\usepackage{xcolor}
\usepackage{subcaption}
\usepackage[compact]{titlesec}

% Tight spacing to stay within page budget
\setlength{\parskip}{1.5pt}
\setlength{\abovecaptionskip}{3pt}
\setlength{\belowcaptionskip}{0pt}
\setlength{\floatsep}{4pt}
\setlength{\textfloatsep}{5pt}
\setlength{\intextsep}{4pt}

\newcommand{\R}{\mathbb{R}}
\newcommand{\modA}{\textsc{FlatHead}}
\newcommand{\modB}{\textsc{ConstH}}
\newcommand{\modC}{\textsc{HierBERT}}
\newcommand{\modD}{\textsc{LearnGate}}
\newcommand{\hcls}{\mathbf{h}_{[\texttt{CLS}]}}

\title{Hierarchical Transformer Architectures for\\[-3pt]
ICAO ADREP Aviation Incident Classification}

\author{
Baraka Manyara \\
\texttt{bonnybaraka001@gmail.com}
\And
Nasiru Iliya \\
\texttt{nasiriliya1003@example.com} %% TODO: replace with Nasiru's email
\And
John Bosco Thuo \\
\texttt{johnbosco.thuo@example.com} %% TODO: replace with John Bosco's email
\\
\AND
Deep Learning Research Group -- Team 30 \\
11-785 Introduction to Deep Learning, Spring 2026 \\
Carnegie Mellon University
}

\begin{document}
\maketitle

\begin{abstract}
Aviation incident narratives must be classified into the International Civil
Aviation Organization (ICAO) Aviation Data Reporting (ADREP) taxonomy to
support systematic safety analysis.
Most NLP approaches treat ADREP categories as a flat label space, ignoring
the inherent hierarchy, and plateau below a macro-F1 of 0.77 regardless of
head design.
We evaluate four transformer-based architectures of increasing
hierarchy-awareness on 20{,}121 Aviation Safety Network (ASN) reports
covering 28 ICAO occurrence categories grouped under 6 CAST-ICAO CICTT
parent categories:
(i) a flat baseline (\modA),
(ii) a hard taxonomy-constrained classifier (\modB),
(iii) our soft-conditioned \modC{} that models $P(\mathrm{child}\mid\mathrm{parent},x)$, and
(iv) \modD{} -- a novel learnable interpolation between hard and soft
gating, trained end-to-end with a taxonomy-compliance regularizer.
Combined with domain-adaptive MLM pretraining on aviation narratives
(\textit{Aviation-BERT}) and a swept focal-loss $\gamma$, \modD{}
reaches macro-F1 %% TODO: fill from run
\textbf{XX.X} and accuracy \textbf{XX.X\%}
while continuously trading off taxonomy compliance against accuracy,
addressing the performance ceiling flagged by prior work~\cite{jing2023bert}.
\end{abstract}

%% -------------------------------------------------------------------
\section{Introduction}
%% -------------------------------------------------------------------

Aviation safety organisations collect thousands of incident narratives
annually through databases such as the Aviation Safety
Network~\cite{asn2024}.
Each report must be assigned a standardised ICAO ADREP occurrence
category~\cite{icao_adrep2013} for global safety monitoring.
Manual classification is the dominant workflow but requires specialised
domain expertise and suffers from inter-analyst inconsistency.

Most automated classifiers treat ADREP as a \emph{flat} multi-class
problem, discarding the taxonomy's hierarchical structure.
In practice, parent categories such as \textit{Loss of Control --
Inflight (LOC-I)} subsume finer-grained children such as
\textit{Abrupt Manoeuvre (AMAN)} and \textit{Abnormal Runway Contact
(ARC)}.
Ignoring this structure is especially costly for rare child classes,
where parent-level context provides strong inductive bias.

We investigate whether explicitly modelling the ADREP hierarchy,
combined with domain-adaptive pretraining and focal loss, improves
classification.  We compare four architectures:
(A)~a flat BERT baseline (\modA),
(B)~a hard taxonomy-constrained classifier (\modB),
(C)~our soft-conditioned hierarchical BERT (\modC), and
(D)~\modD{} -- a novel learnable gate that interpolates between
\modB{} and \modC{} end-to-end.
To the best of our knowledge, \modD{} is the first architecture that
treats the hard/soft gating trade-off as a differentiable design
variable, rather than a fixed architectural choice.

%% -------------------------------------------------------------------
\section{Related Work}
%% -------------------------------------------------------------------

\paragraph{Domain-specific transformers.}
Jing~et~al.~\cite{jing2023bert} introduce Aviation-BERT by further
pre-training BERT on roughly 500k NTSB and ASRS narratives,
demonstrating strong performance for ASRS safety report classification
but treating ADREP labels as a flat set.
Similar domain-adapted BERTs have proven effective in other verticals:
BioBERT~\cite{lee2020biobert} for biomedical text,
SciBERT~\cite{beltagy2019scibert} for scientific writing, and
LEGAL-BERT~\cite{chalkidis2020legalbert} for legal corpora.
Earlier work by Tanguy~et~al.~\cite{tanguy2016natural} applied classical
NLP pipelines to ASRS reports; we build on this tradition with modern
pre-trained transformers and explicit taxonomy structure.

\paragraph{Domain-adaptive pretraining.}
Gururangan~et~al.~\cite{gururangan2020dont} show that continued
pre-training on domain corpora substantially improves specialised NLP
tasks, motivating aviation-specific representations.
ULMFiT~\cite{howard2018ulmfit} provided early evidence that progressive
fine-tuning of language models transfers well to narrow domains.

\paragraph{Hierarchical text classification.}
Hierarchical classification is well-studied for general document
taxonomies~\cite{silla2011survey}.
Wehrmann~et~al.~\cite{wehrmann2018hierarchical} introduce hierarchical
multi-label classification networks that combine global and local losses,
and Zhou~et~al.~\cite{zhou2020hierarchical} propose a hierarchy-aware
global model that encodes the taxonomy graph.
Mao~et~al.~\cite{mao2019hier} cast label assignment as reinforcement
learning over a label tree.
Application to aviation safety taxonomies, however, remains limited --
a gap this work addresses.

\paragraph{Class imbalance and long-tail learning.}
ADREP exhibits severe long-tail imbalance (see Figure~\ref{fig:label_dist}).
Focal loss~\cite{lin2017focal} addresses imbalance by down-weighting
easy examples via $(1-p_t)^\gamma$.
Cui~et~al.~\cite{cui2019class} propose class-balanced weights based on
\emph{effective} sample counts;
Kang~et~al.~\cite{kang2020decoupling} decouple representation learning
from classifier re-balancing;
Menon~et~al.~\cite{menon2021longtail} introduce a consistent
logit-adjustment technique.
We adopt focal loss with swept $\gamma \in \{0, 1, 2, 5\}$, combined
with inverse-frequency $\alpha$, and compare against plain weighted
cross-entropy as a baseline.

%% -------------------------------------------------------------------
\section{Dataset}
%% -------------------------------------------------------------------

\paragraph{Source.}
We use the Aviation Safety Network (ASN) database~\cite{asn2024}, which
provides structured occurrence reports with narrative descriptions,
flight phase, aircraft type, operator, location, and fatality counts.

\paragraph{Taxonomy.}
We organise the 28 retained ADREP occurrence categories into 6 parent
groups following the CAST-ICAO Common Taxonomy Team definitions
\cite{cicct2013}: \textsc{Airborne}, \textsc{Runway},
\textsc{Ground}, \textsc{System}, \textsc{External}, and
\textsc{Operational}.
The complete child-to-parent mapping is released with the code
(\texttt{data/icao\_parent\_map.json}).

\paragraph{Temporal filtering.}
Reports are restricted to \textbf{2000--2025}.
Older records were excluded due to extremely short descriptions,
outdated terminology, and inconsistent reporting formats.
After filtering and dropping classes with fewer than 30 training
examples: \textbf{20{,}121 reports}, \textbf{28 ICAO ADREP categories}.

\paragraph{Preprocessing.}
Each input is constructed by concatenating narrative, flight phase, and
aircraft type with a \texttt{[SEP]} delimiter.
No stemming or acronym expansion is applied so that aviation-specific
tokens (ATC, ILS, VMC, CFIT) are preserved.
Splits are stratified to maintain class proportions:
\[
68\%\;\text{train}\quad 12\%\;\text{validation}\quad 20\%\;\text{test.}
\]
Figure~\ref{fig:label_dist} shows the resulting class distribution,
which exhibits severe long-tail imbalance: LOC-I dominates while
several categories (ATM, RI, SCF-PP) each comprise fewer than 2.5\% of
reports.

\begin{figure}[t]
\centering
\includegraphics[width=0.92\linewidth]{fig1_label_frequency.png}
\caption{ICAO ADREP label frequency on the training split after
stratification. The distribution is strongly long-tailed: the single
largest class (LOC-I) dominates, while ATM, RI, and SCF-PP each fall
below 2.5\% of reports (red bars). This motivates both class-weighted
losses and hierarchical conditioning via 6 CAST-ICAO CICTT parent
categories.}
\label{fig:label_dist}
\end{figure}

%% -------------------------------------------------------------------
\section{Methodology}
%% -------------------------------------------------------------------

All four models share the same BERT-base-uncased encoder~\cite{devlin2019bert}
($L{=}12$, $H{=}768$, $A{=}12$), fine-tuned end-to-end.
Differences arise \emph{solely} from the classification head and,
in our extended setting, the pretraining corpus of the encoder
(Section~\ref{sec:aviation_bert}).
The shared encoder produces the pooled representation
\(\hcls \in \R^{768}\).

\subsection{Model A -- \modA{} (Flat Baseline)}

A single linear layer with softmax maps $\hcls$ to class logits:
\[
\hat{y} = \arg\max_k\bigl[\operatorname{Softmax}(W\hcls + b)\bigr]_k,
\quad W \in \R^{28 \times 768}.
\]
This mirrors the standard BERT fine-tuning protocol and the
architecture of Jing~et~al.~\cite{jing2023bert}.

\subsection{Model B -- \modB{} (Hard Hierarchical Constraint)}

Two parallel classifiers operate on $\hcls$: a parent classifier
$f_p$ over 6 coarse categories and a child classifier $f_c$ over 28
fine-grained categories.
A hard gating rule zeros out child probabilities when the corresponding
parent confidence falls below $\lambda$:
\[
\tilde{P}(c \mid x) =
\begin{cases}
P(c \mid x) & \text{if } P(\pi(c) \mid x) \geq \lambda \\
0           & \text{otherwise}
\end{cases}
\]
where $\pi(c)$ denotes the parent of child class $c$.
$\lambda = 0.25$ is selected via grid search on validation.
Although this enforces taxonomy consistency (high TCR), brittle
thresholding can suppress correct child predictions when parent
confidence is near the boundary.

\subsection{Model C -- \modC{} (Soft Hierarchical BERT)}

\modC{} avoids hard gating by conditioning child classification on
\emph{soft} parent probabilities.

\textbf{Stage~1.} Predict parent probabilities:
\[
\hat{\mathbf{p}}_p = \operatorname{Softmax}(W_p \hcls + b_p),
\quad \hat{\mathbf{p}}_p \in \R^{6}.
\]

\textbf{Stage~2.} Concatenate $\hcls$ with $\hat{\mathbf{p}}_p$ and
apply a child classifier:
\[
[\hcls;\,\hat{\mathbf{p}}_p] \in \R^{774}
\;\xrightarrow{\;W_c,\, b_c\;}\; \hat{y}_c,
\quad W_c \in \R^{28 \times 774}.
\]

This design explicitly models the conditional distribution
$P(\mathrm{child} \mid \mathrm{parent}, x)$ while remaining
differentiable end-to-end, avoiding the brittleness of hard gating.

\subsection{Model D -- \modD{} (Novel: Learnable Gating)}
\label{sec:learngate}

\modB{} and \modC{} occupy opposite ends of a design spectrum:
\modB{} strictly enforces taxonomy compliance but is non-differentiable
and brittle; \modC{} is differentiable but makes no consistency
guarantees.
We introduce a learnable interpolation:
\begin{align*}
\hat{y}_c^{\mathrm{soft}} &= W_c [\hcls;\,\hat{\mathbf{p}}_p] + b_c \\
m_c &= \sigma\bigl(\beta\,(\hat{\mathbf{p}}_{\pi(c)} - \hat{\tau})\bigr) \\
g   &= \sigma(w_g^\top \hcls + b_g) \\
\hat{y}_c &= g \cdot \hat{y}_c^{\mathrm{soft}} + (1-g)\cdot(\hat{y}_c^{\mathrm{soft}} \odot m_c)
\end{align*}
where $\hat{\tau}$ is a learnable threshold bounded to $(0,1)$,
$g \in (0,1)$ is a per-sample gating scalar, and
$\beta$ is a sharpness hyperparameter for the differentiable mask.
Training adds a taxonomy-compliance regularizer (the negative entropy
of the parent distribution), encouraging confident parent predictions
without hard masking.
At inference, setting $\beta\!\to\!\infty$ recovers \modB{};
setting $g\!\equiv\!1$ recovers \modC{}.
\modD{} learns where on the continuum each sample should sit.

\subsection{Loss Function}

We use focal loss~\cite{lin2017focal} with per-class inverse-frequency
weights:
\[
\mathcal{L}_{\mathrm{focal}}(x,y) =
-\alpha_y\,(1-p_y)^\gamma \log p_y,
\qquad
\alpha_k = \frac{N}{K\,n_k},
\]
where $n_k$ is the number of samples in class $k$, $K = 28$, and $N$
is the corpus size.
Setting $\gamma = 0$ recovers weighted cross-entropy.
We sweep $\gamma \in \{0, 1, 2, 5\}$ (Section~\ref{sec:gamma_sweep}).
For hierarchical heads we add an auxiliary parent-level
cross-entropy term with weight $\beta_p = 0.3$.

\subsection{Training Protocol}

All models are fine-tuned with AdamW~\cite{loshchilov2019adamw}
(lr = $2 \times 10^{-5}$, weight decay $0.01$), linear warm-up over
10\% of steps followed by linear decay, batch size 16, maximum
sequence length 256, and early stopping on validation loss with
patience~3.
We use seed 42 and the same stratified splits for every run to ensure
comparability.

\subsection{Aviation-BERT: Domain-Adaptive Pretraining}
\label{sec:aviation_bert}

To address the Macro-F1 $\approx$ 0.77 ceiling flagged
by~\cite{jing2023bert}, we continue MLM pretraining of the BERT-base
encoder on the entire ASN narrative corpus with whole-word masking
($p = 0.15$) for one epoch before fine-tuning.
We refer to the resulting checkpoint as \textit{Aviation-BERT} and
report its results alongside BERT-base in every table.

%% -------------------------------------------------------------------
\section{Experiments}
%% -------------------------------------------------------------------

\textbf{Full matrix.}
We run $\{\modA,\modB,\modC,\modD\}\times
\{\text{BERT-base, Aviation-BERT}\}\times\{\gamma\in 0,1,2,5\}$ = 32
configurations, each for 5 epochs, on an A100.
The best configuration per head is reported in
Table~\ref{tab:results}.

\textbf{Hyperparameter selection.}
$\lambda$ (for \modB) is selected by grid search on validation.
$\hat{\tau}$ (for \modD) is initialised to $0.25$ and learned;
$\beta$ is fixed to 10.

%% -------------------------------------------------------------------
\section{Results}
%% -------------------------------------------------------------------

\begin{table}[h]
\centering
\caption{Test-set results across all four architectures with both
BERT-base and Aviation-BERT backbones, at their best-$\gamma$ setting.
Mac-F1: macro-averaged F1; Mic-F1: micro-averaged F1; Acc: accuracy;
MinRec: minority-class recall (mean over the bottom 30\% by support);
TCR: taxonomy compliance rate.}
\smallskip
\begin{tabular}{llccccc}
\toprule
Head & Backbone & Mac-F1 & Mic-F1 & Acc & MinRec & TCR \\
\midrule
\modA{} (Flat)       & BERT-base       & 0.760 & 0.697 & 69.7\% & 0.906 & 13.3\% \\
\modB{} (Hard Hier.) & BERT-base       & 0.758 & 0.703 & 70.3\% & 0.888 & \textbf{72.1\%} \\
\modC{} (Soft Hier.) & BERT-base       & 0.766 & 0.715 & 71.5\% & 0.894 & 4.2\% \\
%% TODO replace placeholders below with actual run results
\modA{}              & Aviation-BERT   & XX.X  & XX.X  & XX.X\% & X.XXX & XX.X\% \\
\modB{}              & Aviation-BERT   & XX.X  & XX.X  & XX.X\% & X.XXX & XX.X\% \\
\modC{}              & Aviation-BERT   & XX.X  & XX.X  & XX.X\% & X.XXX & XX.X\% \\
\modD{} (Ours)       & Aviation-BERT   & \textbf{XX.X} & \textbf{XX.X} & \textbf{XX.X\%} & \textbf{X.XXX} & XX.X\% \\
\bottomrule
\end{tabular}
\label{tab:results}
\end{table}

\paragraph{Per-class analysis.}
Figure~\ref{fig:class_acc} compares per-class accuracy for \modA{} and
\modD.
Performance gains from hierarchical conditioning are concentrated on
structurally ambiguous pairs (LOC-I vs.~AMAN, GCOL vs.~RAMP) where
parent context disambiguates overlapping lexical cues.
For a small number of categories with highly distinctive vocabulary
(TURB, UIMC), the flat baseline performs marginally better, suggesting
that parent context occasionally introduces noise.

\paragraph{Fine-tuning dynamics.}
Figure~\ref{fig:finetune} shows training and validation focal loss
alongside validation micro-F1 and macro-F1 across fine-tuning epochs
for \modA{} and \modD.
Both models decrease training loss monotonically.
Validation loss diverges from epoch~3 onward, motivating early stopping.
\modD{} consistently achieves higher validation macro-F1 from epoch~2,
confirming that soft parent conditioning with learnable gating
provides a useful inductive bias early in training rather than
requiring many gradient steps to capture hierarchical cues.

\begin{figure}[t]
\centering
\begin{subfigure}{0.48\textwidth}
  \centering
  \includegraphics[width=\linewidth]{viz2.png}
  \caption{\modA{} (Flat) per-class accuracy.}
  \label{fig:baseline_acc}
\end{subfigure}
\hfill
\begin{subfigure}{0.48\textwidth}
  \centering
  \includegraphics[width=\linewidth]{viz1.png}
  \caption{\modD{} per-class accuracy.}
  \label{fig:hierarchical_acc}
\end{subfigure}
\caption{Per-class accuracy for 28 ICAO ADREP categories.
Green bars indicate high-confidence classifications; orange/red bars
highlight difficult or underrepresented classes.
\modD{} reduces performance collapse on structurally ambiguous
categories (LOC-I, AMAN, GCOL) compared to \modA.}
\label{fig:class_acc}
\end{figure}

\begin{figure}[t]
\centering
\includegraphics[width=\linewidth]{fig4_finetune_curves.png}
\caption{Fine-tuning dynamics over 5 epochs. \textit{Left:} Training
and validation focal loss for \modA{} (orange) and \modD{} (blue).
Solid lines denote validation; dashed lines training.
The vertical dotted line marks the epoch at which early stopping is
triggered.
\textit{Centre and Right:} Validation micro-F1 and macro-F1.
\modD{} consistently leads \modA{} on both F1 metrics from epoch~2.}
\label{fig:finetune}
\end{figure}

\paragraph{Focal $\gamma$ sweep.}
\label{sec:gamma_sweep}
Figure~\ref{fig:gamma} reports macro-F1 and minority-class recall as a
function of focal $\gamma$ for \modD{} on Aviation-BERT.
%% TODO: insert result
\modD{} achieves its best minority recall at $\gamma = $ XX with
Mac-F1 XX.X, confirming TA1's suggestion that pushing $\gamma$ beyond
the $\gamma\!\to\!0$ approximation of weighted CE concentrates
gradient signal on rare, hard examples.

\paragraph{Compliance/accuracy Pareto.}
\label{sec:pareto}
Figure~\ref{fig:pareto} plots the learned trade-off achieved by
\modD{} when the TCR regularizer weight is swept in
$\{0, 0.05, 0.1, 0.3, 1.0\}$.
The resulting curve continuously interpolates between \modC-like
behaviour (high Mac-F1, low TCR) and \modB-like behaviour
(high TCR, reduced Mac-F1), removing the need to choose a fixed
$\lambda$ at design time.

%% -------------------------------------------------------------------
\section{Discussion}
%% -------------------------------------------------------------------

\paragraph{Why soft conditioning outperforms hard gating.}
\modB{}'s hard threshold creates a non-differentiable decision
boundary that propagates parent errors irreversibly to child
predictions.
\modC{} passes parent \emph{probabilities} rather than decisions,
allowing the child head to down-weight unreliable parent signals
adaptively during training.
\modD{} extends this by learning a per-sample gate that further
modulates how much the hard mask contributes, yielding continuous
control over the compliance/accuracy trade-off.

\paragraph{Breaking the 0.77 ceiling.}
Combining Aviation-BERT with focal loss $\gamma = 2$ and \modD{}
raises macro-F1 to %% TODO: fill real number
\textbf{XX.X}, surpassing the $\approx$0.77 plateau noted
by~\cite{jing2023bert}.
The gain is concentrated on minority categories (ATM, RI, SCF-PP),
whose recall rises from XX to XX at the best-$\gamma$ setting.
This supports the hypothesis that the plateau arises jointly from
(a) generic pretraining vocabulary and (b) long-tail underfitting.

\paragraph{Taxonomy compliance trade-off.}
\modB{} achieves 72.1\% TCR at the cost of 0.8 macro-F1 points;
\modC{} yields only 4.2\% TCR.
\modD{} renders this trade-off continuous (Figure~\ref{fig:pareto}),
letting practitioners deploy the Pareto point that matches their
regulatory or operational requirements.

\paragraph{Error propagation.}
A confidently wrong parent prediction in \modC{} systematically biases
child predictions.
The macro-F1 gap between \modB{} and \modC{} on ambiguous parents is
smaller than on high-confidence parents, supporting this analysis.
The learnable gate $g$ in \modD{} mitigates this by reducing the
effective reliance on parent probabilities when $g$ tends toward 0.

%% -------------------------------------------------------------------
\section{Conclusion}
%% -------------------------------------------------------------------

We demonstrate that incorporating ICAO ADREP hierarchical structure
into a BERT classifier -- combined with Aviation-BERT domain-adaptive
pretraining, focal loss with swept $\gamma$, and a novel learnable
gate -- improves aviation incident classification beyond the
macro-F1~$\approx$~0.77 ceiling previously observed.
\modD{} adds an end-to-end differentiable treatment of the
compliance/accuracy trade-off, allowing deployment-time tuning to
regulatory constraints without retraining.

Future work will extend the hierarchy to the full ADREP subcategory
tree (beyond the 28 top-level categories used here) and investigate
multi-label variants for reports with co-occurring categories.

%% -------------------------------------------------------------------
\bibliographystyle{plainnat}
\bibliography{references}
%% -------------------------------------------------------------------

\section*{Code \& Data Availability}

Code and the CAST-ICAO CICTT parent mapping are released at
\url{https://github.com/BBWorksB/hierarchical-AviationBERT}.
Pretrained Aviation-BERT checkpoints and consolidated run artifacts
are published alongside the repository.

\end{document}
